\documentclass{article}
\usepackage[utf8]{inputenc}
\usepackage[english]{babel}

\begin{document}

\section*{Reformed Anthropology and Eschatological Dichotomy: Key Components with Greek Terms}

1. \textbf{Flesh (\textit{sarx})} \\
   \hspace*{0.5cm} Refers to the sinful, fallen nature of humanity under sin (e.g., Romans 8:3, Galatians 5:16-17). In the "old age," \textit{sarx} signifies life in bondage to sin, a mode of existence characterized by rebellion against God and inability to please Him (Romans 8:7-8). Contrastingly, in the "new creation," the dominion of \textit{sarx} is overcome through the Spirit.

2. \textbf{Body (\textit{sōma})} \\
   \hspace*{0.5cm} Refers to the physical body but also carries theological significance. In the "old man," the \textit{sōma} is presented as an instrument for sin (Romans 6:12). In the "new man," the body becomes an instrument for righteousness, destined for transformation and resurrection (Romans 12:1; 1 Corinthians 15:42-44).

3. \textbf{Mind (\textit{nous})} \\
   \hspace*{0.5cm} Represents intellect, reasoning, and understanding. The \textit{nous} in the "old age" is hostile to God and futile in its thinking (Romans 1:28). In the "new creation," the mind is renewed, enabling believers to discern and live according to God's will (Romans 12:2; Ephesians 4:23).

4. \textbf{Soul (\textit{psuchē})} \\
   \hspace*{0.5cm} Refers to the whole person, encompassing emotions and desires. In the "old man," the \textit{psuchē} is bound to the "present evil age" and its corrupted desires. In the "new man," it points to eternal life and existence in union with Christ (Matthew 16:26, Philippians 1:27).

5. \textbf{Heart (\textit{kardia})} \\
   \hspace*{0.5cm} Symbolizes the center of emotions, will, and spiritual life. In the "old man," the \textit{kardia} is hardened and deceitful (Jeremiah 17:9). In the "new man," it is the locus of transformation, where the law of God is written and faith arises (Romans 10:10; 2 Corinthians 3:3).

6. \textbf{Spirit (\textit{pneuma})} \\
   \hspace*{0.5cm} Refers to the immaterial aspect of humanity that relates to God. In the "old age," the human \textit{pneuma} is lifeless without the Holy Spirit. In the "new creation," it is united with the Holy Spirit for transformation and empowerment to live a God-pleasing life (Romans 8:16; 1 Corinthians 2:12).

7. \textbf{Will (\textit{thelēma})} \\
   \hspace*{0.5cm} Refers to human volition and the ability to choose. In the "old man," the \textit{thelēma} is enslaved to sin, unable to do the good it desires (Romans 7:18-23). In the "new man," the will is liberated to obey God and seek His purposes.

8. \textbf{Conscience (\textit{suneidēsis})} \\
   \hspace*{0.5cm} Represents the internal moral compass. In the "old age," the conscience may be weak, seared, or defiled (1 Timothy 4:2; 1 Corinthians 8:7). In the "new creation," it is purified and serves as a guide for living in righteousness (Hebrews 9:14).

\section*{Eschatological Terms in Paul's Dichotomy}

9. \textbf{Cosmos (\textit{kosmos})} \\
   \hspace*{0.5cm} Refers to the world as a system opposed to God. It represents the domain of sin and death, the "old age" under the dominion of the flesh (Galatians 1:4). In contrast, the "new creation" transcends the \textit{kosmos}, reflecting a redeemed order under Christ's lordship.

10. \textbf{Age (\textit{aiōn})} \\
   \hspace*{0.5cm} Contrasts the "present evil age" (\textit{aiōn houtos}, Galatians 1:4), marked by sin and death, with the "age to come," inaugurated by Christ's resurrection. The transition from one age to another defines the believer's eschatological hope (Ephesians 1:21).

11. \textbf{Old Man (\textit{palaios anthrōpos}) and New Man (\textit{kainos anthrōpos})} \\
   \hspace*{0.5cm} The "old man" signifies the self in Adam, enslaved to sin and identified with the "present evil age" (Romans 6:6). The "new man" refers to the regenerate self, created in righteousness and aligned with the "new creation" (Ephesians 4:22-24).

12. \textbf{New Creation (\textit{kainē ktisis})} \\
   \hspace*{0.5cm} Represents the eschatological reality inaugurated by Christ's redemptive work (2 Corinthians 5:17). Believers participate in the \textit{kainē ktisis}, experiencing freedom from sin's power and the restoration of all things.

13. \textbf{Bondage to Sin and Freedom in Christ} \\
   \hspace*{0.5cm} Bondage to sin characterizes life in the "old age," where humanity is enslaved to sin and death (Romans 6:16). Freedom in Christ marks the "new creation," where believers are liberated to live in obedience and righteousness (Romans 6:18).

14. \textbf{Walking in the Flesh vs. Walking in the Spirit} \\
   \hspace*{0.5cm} Walking in the flesh (\textit{sarx}) reflects the lifestyle of the "old man," dominated by sin (Galatians 5:19-21). Walking in the Spirit (\textit{pneuma}) is the hallmark of the "new man," empowered by the Holy Spirit for godly living (Galatians 5:22-25).

15. \textbf{Old Age (\textit{palaios aiōn}) and New Age (\textit{kainos aiōn})} \\
   \hspace*{0.5cm} The "old age" refers to the era of sin and death under Adam, characterized by the dominion of the flesh. The "new age" is the eschatological reality inaugurated by Christ, marked by life, righteousness, and the Spirit (Romans 8:2).

\end{document}
